\documentclass[]{article}

\usepackage{amsmath}
\usepackage{multirow}
\usepackage{natbib}
\usepackage{graphicx} 


\begin{document}

Our analysis proceeds in several stages, beginning with the identification of distinct employee trajectories using Latent Class Analysis (LCA), followed by an examination of the organizational and psychological factors that predict membership in these classes.

\subsection{Latent class trajectories of AI-induced job security}
To identify unobserved heterogeneity in employees' perceptions of AI's impact on their job security, we conducted an LCA on the joint distribution of current and expected job security. We evaluated models with two to six latent classes, using the Bayesian Information Criterion (BIC) and normalized entropy to determine the optimal solution. As shown in Figure \ref{fig:lca_selection}, the three-class solution provided the optimal fit to the data. This model yielded the lowest BIC value (14,592.21) and a strong normalized entropy of 0.8032, indicating a parsimonious fit with excellent class separation and classification certainty.

\begin{figure}[h!]
    \centering
    \includegraphics[width=0.6\textwidth]{../input_files/plots/step_2_lca_bic_entropy_1776114673.png}
    \caption{Model selection metrics for the Latent Class Analysis (LCA) of AI-induced job security perceptions. The plot displays the Bayesian Information Criterion (BIC; blue line) and normalized entropy (red line) for models with two to six latent classes. The three-class solution was selected as the optimal model, as it exhibits the lowest BIC value and a high normalized entropy, indicating the best balance of model fit and parsimony with clear class separation.}
    \label{fig:lca_selection}
\end{figure}

The three emergent classes represent distinct psychological trajectories:
\begin{itemize}
    \item \textbf{Class 1: "Resiliently Optimistic" (60.9\% of the sample):} This majority class is characterized by a highly positive outlook. Members of this group overwhelmingly report that AI has a positive impact on their current job security and expect this trend to continue.
    \item \textbf{Class 2: "Stagnant Neutral" (26.5\% of the sample):} This class exhibits a strong anchoring to the status quo. These employees report that AI has had no impact on their job security and largely expect this to remain the case.
    \item \textbf{Class 3: "Anxiously Declining" (12.6\% of the sample):} This minority class perceives AI as a direct threat. Members report a negative current impact on their job security and anticipate that this negative impact will persist or worsen.
\end{itemize}

The conditional response probabilities for each class, detailed in Table \ref{tab:lca_probs}, provide a granular view of these trajectories. For instance, individuals in the "Resiliently Optimistic" class have a cumulative probability of 84.2\% of reporting a positive current impact, while those in the "Anxiously Declining" class have a 92.4\% probability of reporting a negative current impact.

\begin{table*}[t]
\centering
\caption{LCA Class Sizes and Conditional Response Probabilities for the 3-Class Solution.}
\label{tab:lca_probs}
\resizebox{\textwidth}{!}{
\begin{tabular}{l|cccccc|cccccc}
\hline
\hline
& \multicolumn{6}{c|}{\textbf{Current Impact on Job Security}} & \multicolumn{6}{c}{\textbf{Expected Impact on Job Security (3 years)}} \\
\textbf{Latent Class} & Sig. Neg. & Sl. Neg. & No Impact & Sl. Pos. & Sig. Pos. & Not Sure & Sig. Neg. & Sl. Neg. & No Impact & Sl. Pos. & Sig. Pos. & Not Sure \\
\hline
\textbf{1. Resiliently Optimistic} & 0.000 & 0.030 & 0.124 & 0.419 & 0.423 & 0.004 & 0.000 & 0.046 & 0.120 & 0.377 & 0.449 & 0.008 \\
(60.90\%) & & & & & & & & & & & & \\
\textbf{2. Stagnant Neutral} & 0.000 & 0.051 & 0.861 & 0.052 & 0.000 & 0.037 & 0.000 & 0.152 & 0.628 & 0.173 & 0.000 & 0.047 \\
(26.51\%) & & & & & & & & & & & & \\
\textbf{3. Anxiously Declining} & 0.222 & 0.701 & 0.041 & 0.036 & 0.000 & 0.000 & 0.296 & 0.510 & 0.130 & 0.064 & 0.000 & 0.001 \\
(12.59\%) & & & & & & & & & & & & \\
\hline
\end{tabular}
}
\end{table*}

\subsection{Dimensionality reduction of affective dispositions}
To model employees' underlying emotional attitudes towards AI without introducing multicollinearity in subsequent models, we performed an Exploratory Factor Analysis (EFA) on 16 binary items. The analysis, using a tetrachoric correlation matrix and Varimax rotation, successfully extracted two distinct latent factors. As visualized in Figure \ref{fig:factor_loadings}, these factors clearly represent "Positive Affect" (e.g., 'optimistic', 'empowered', 'motivated') and "Negative Affect" (e.g., 'worried', 'threatened', 'anxious'). Both factors demonstrated high internal consistency, with McDonald's Omega values of 0.9004 for Positive Affect and 0.9011 for Negative Affect.

Further validation confirmed that the 'Negative Affect' factor was significantly associated with key deterrents to AI adoption, showing correlations with distrust of AI accuracy ($r = -0.246, p < 0.001$), fear of job loss ($r = -0.173, p < 0.001$), and privacy concerns ($r = -0.127, p < 0.001$). This confirms that the latent factor effectively captures underlying psychological frictions related to AI integration. The resulting factor scores were used as continuous covariates in subsequent models.

\begin{figure}[h!]
    \centering
    \includegraphics[width=0.6\textwidth]{../input_files/plots/step_3_qhd_factor_loadings_1776114809.png}
    \caption{Heatmap of factor loadings from an Exploratory Factor Analysis on 16 emotional attitude items after a Varimax rotation. The analysis reveals two distinct latent factors. The first factor, "Positive Affect," is strongly loaded by items such as 'Optimistic' and 'Motivated'. The second factor, "Negative Affect," is characterized by items like 'Uncertain', 'Skeptical', and 'Anxious', confirming the successful reduction of emotional attitudes into two coherent constructs.}
    \label{fig:factor_loadings}
\end{figure}

\subsection{Identifying key organizational predictors with elastic net}
We next employed an Elastic Net logistic regression to identify the most salient organizational policies and cultural attributes predicting membership in the "Resiliently Optimistic" class. This method performs automated feature selection, isolating the most impactful predictors from a larger set of correlated variables.

The results, shown in Figure \ref{fig:elastic_net}, reveal that organizational catalysts for optimism are concentrated around incentive structures and employee agency. The strongest positive predictors were non-monetary incentives, such as peer recognition ($\beta = 0.303$) and learning benefits or certifications ($\beta = 0.236$). Structural enablers that provide employees with agency, such as direct involvement in the AI development process ($\beta = 0.129$) and the freedom to choose AI tools ($\beta = 0.124$), were also highly significant. Conversely, the strongest negative predictors were deterrents rooted in risk and uncertainty, namely distrust in AI accuracy ($\beta = -0.158$) and privacy concerns ($\beta = -0.141$). Notably, the stated importance of regular training yielded a coefficient of zero, suggesting that abstract commitments to training, without tangible incentives or participatory frameworks, are insufficient to foster optimism.

\begin{figure}[h!]
    \centering
    \includegraphics[width=0.6\textwidth]{../input_files/plots/step_4_elastic_net_coefficients_1776115025.png}
    \caption{Elastic Net logistic regression coefficients identifying predictors of membership in the "Resiliently Optimistic" class. Non-monetary incentives (e.g., peer recognition, learning benefits) and structural enablers (e.g., employee involvement) are the most potent positive predictors. Conversely, deterrents such as distrust in AI accuracy and privacy concerns are the strongest negative predictors.}
    \label{fig:elastic_net}
\end{figure}

\subsection{Multinomial logistic regression of class membership}
Building on the feature selection, we estimated a multinomial logistic regression model to evaluate the adjusted odds of membership in the "Resiliently Optimistic" or "Stagnant Neutral" classes relative to the "Anxiously Declining" reference class. The model included key organizational enablers, deterrents, and their interactions.

The odds ratios, visualized in Figure \ref{fig:odds_ratios}, confirm and extend the Elastic Net findings. Organizational enablers significantly increased the odds of belonging to the "Resiliently Optimistic" class. For example, providing career progression incentives increased the odds by 15.6\% (OR = 1.156), while direct employee involvement in AI development increased them by 11.1\% (OR = 1.111). The most powerful determinant was the fear of job loss; employees who cited this as a deterrent had their odds of being in the optimistic class reduced by over 82\% (OR = 0.174) compared to those who did not.

The model also revealed significant interaction effects. First, the freedom to choose AI tools was found to buffer the negative impact of privacy concerns; granting employees autonomy over tool selection significantly weakened the negative effect of privacy anxieties on their odds of being optimistic (OR for interaction = 1.281). Conversely, a second interaction showed that emphasizing the importance of regular training in a context where employees fear accountability for AI-related mistakes may exacerbate anxiety. This combination significantly reduced the odds of belonging to the optimistic class (OR for interaction = 0.712), suggesting that training initiatives must be paired with supportive accountability frameworks.

\begin{figure}[h!]
    \centering
    \includegraphics[width=0.6\textwidth]{../input_files/plots/step_5_forest_plot_odds_ratios_1776115179.png}
    \caption{Forest plot of odds ratios from the multinomial logistic regression, showing predictors of membership in the "Resiliently Optimistic" class relative to the "Anxiously Declining" class. Deterrents like worry about job loss (QF\_3) drastically decrease the odds of an optimistic outlook, while enablers like career progression incentives (QKC\_4) and employee involvement (QKB\_2\_11) increase them. Interaction effects show that autonomy over tools (QKB\_2\_9\_x\_QF\_7) can buffer the negative impact of privacy anxieties.}
    \label{fig:odds_ratios}
\end{figure}

\subsection{Simulating the impact of organizational policies}
To translate these regression findings into actionable insights, we conducted a Monte Carlo policy simulation to estimate the predicted probability of an average employee belonging to the "Resiliently Optimistic" class under different organizational scenarios. We specifically contrasted the effects of providing regular training with direct employee involvement in AI development.

The results, shown in Figure \ref{fig:policy_sim}, demonstrate that policies promoting active agency are far more effective than those focused on passive capability building. In the baseline scenario with neither training nor involvement, the predicted probability of being in the optimistic class is approximately 58\%. Introducing training alone provides only a marginal benefit. However, implementing a policy of direct employee involvement in AI development substantially increases this probability to over 65\%. This finding underscores that to foster psychological stability, organizations must grant employees an active, participatory role in the technological transition.

\begin{figure}[h!]
    \centering
    \includegraphics[width=0.6\textwidth]{../input_files/plots/step_6_policy_simulation_probabilities_1776115363.png}
    \caption{Predicted probability of membership in the "Resiliently Optimistic" class from a Monte Carlo policy simulation under four organizational scenarios. The simulation demonstrates that involving employees in the AI development process substantially elevates the probability of a positive outlook. In contrast, providing training alone yields only a marginal improvement, underscoring that active employee participation is a more potent catalyst for psychological stability than passive capability building.}
    \label{fig:policy_sim}
\end{figure}

\subsection{Robustness of findings}
To ensure the validity of our findings, we performed several robustness checks. First, a Variance Inflation Factor (VIF) analysis on the full design matrix of the multinomial model confirmed the absence of severe multicollinearity. All primary predictors exhibited VIF values well below the conventional threshold of 5.0 (e.g., employee involvement in development, VIF = 1.75; fear of accountability, VIF = 2.27).

Second, we conducted a tripartite sensitivity analysis to assess the stability of our results to the handling of "Not sure" responses in the dependent variables. The multinomial logistic regression was re-estimated on three sample variants: (1) the full sample (N=2603), (2) a sample excluding "Not sure" respondents (N=2538), and (3) the full sample with "Not sure" responses imputed to the modal value. The odds ratios for key predictors remained remarkably stable across all three specifications. For instance, the odds ratio for employee involvement (`QKB_2_11`) was highly consistent across the full (OR = 1.111), reduced (OR = 1.119), and imputed (OR = 1.120) samples. Similarly, the powerful negative effect of job loss worry (`QF_3`) was stable (ORs of 0.174, 0.149, and 0.175, respectively). This stability confirms that our findings are robust and not systematically biased by respondent uncertainty.

\end{document}
                